\subsection{Low types, orientation, and the three different centres}

\begin{statement}{A.10}{Known}{The low-type convention table}
Let \(S_{g,n}\) be an oriented closed surface of genus \(g\)
with a set of \(n\) distinct marked points removed.
Write \(\Pi=\pi_1(S_{g,n})\), \(\operatorname{Mod}^{+}\)
for orientation-preserving mapping classes allowing permutations
of those points, \(\operatorname{PMod}^{+}\) for the subgroup
fixing them individually, and \(\operatorname{Mod}^{\pm}\)
for the full group allowing orientation reversal.
All quotients are by \emph{isotopies through homeomorphisms}
preserving the marked set.  The points are punctures, not
boundary circles fixed pointwise.  With \(C_2=\mathbb Z/2\),
the complete table in the six types specified in Rule~5 is
\begin{center}\small
\begin{tabular}{@{}cclll@{}}
\toprule
\((g,n)\)&\(\Pi\)&\(\operatorname{Mod}^{+}\)&
\(\operatorname{PMod}^{+}\)&\(\operatorname{Mod}^{\pm}\)\\
\midrule
\((0,0)\)&\(1\)&\(1\)&\(1\)&\(C_2\)\\
\((0,1)\)&\(1\)&\(1\)&\(1\)&\(C_2\)\\
\((0,2)\)&\(\mathbb Z\)&\(C_2\)&\(1\)&\(C_2\times C_2\)\\
\((1,0)\)&\(\mathbb Z^2\)&\(\operatorname{SL}_2(\mathbb Z)\)&
\(\operatorname{SL}_2(\mathbb Z)\)&\(\operatorname{GL}_2(\mathbb Z)\)\\
\((0,3)\)&\(F_2\)&\(\mathfrak S_3\)&\(1\)&
\(\mathfrak S_3\times C_2\)\\
\((1,1)\)&\(F_2\)&\(\operatorname{SL}_2(\mathbb Z)\)&
\(\operatorname{SL}_2(\mathbb Z)\)&\(\operatorname{GL}_2(\mathbb Z)\)\\
\bottomrule
\end{tabular}
\end{center}
Their centres, keeping the surface group and the two full
mapping-class groups separate, are
\begin{center}\small
\begin{tabular}{@{}ccll@{}}
\toprule
\((g,n)\)&\(Z(\Pi)\)&\(Z(\operatorname{Mod}^{+})\)&
\(Z(\operatorname{Mod}^{\pm})\)\\
\midrule
\((0,0),(0,1)\)&\(1\)&\(1\)&\(C_2\)\\
\((0,2)\)&\(\mathbb Z\)&\(C_2\)&\(C_2\times C_2\)\\
\((1,0)\)&\(\mathbb Z^2\)&\(\{\pm I\}\)&\(\{\pm I\}\)\\
\((0,3)\)&\(1\)&\(1\)&\(\{1\}\times C_2\)\\
\((1,1)\)&\(1\)&\(\{\pm I\}\)&\(\{\pm I\}\)\\
\bottomrule
\end{tabular}
\end{center}
The pure oriented centres are trivial in the sphere rows and
\(\{\pm I\}\) in the torus rows.  The extended pure sphere
groups, and hence their centres, are \(C_2\).
\end{statement}

\begin{proof}
Van Kampen gives \(F_{n-1}\) for a sphere with \(n>0\)
punctures, \(\mathbb Z^2\) for the closed torus, and
\(F_2\) for the once-punctured torus.  These fundamental-group
calculations are independent of a choice of marked-point labels.
The mapping-class computations are verified in
\citet[\S\S2.2.1--2.2.2, Proposition~2.3, Theorem~2.5,
the following once-punctured-torus paragraph, and \S8.1]
{A_low_FarbMargalit2012}; the cited PDF is version~5.0,
pp.~51--57 and 230--231.

Here is the elementary reduction, including the convention-sensitive
parts.  On a sphere with three fixed points, an arc joining two
points in the complement of the third can be isotoped to its
image, with endpoints fixed.  After fixing the arc pointwise,
cutting along it leaves a disc with one marked point.  The
Alexander isotopy, with this point at the centre, fixes both
the point and the boundary.  Thus the pure oriented group is
trivial.  For fewer marked points, extra points can first be
fixed by isotopy in their complement; the same argument applies.
Every permutation of at most three points is realized by an
orientation-preserving homeomorphism.  A reflection of the
sphere fixes three chosen points individually.  Consequently
the homomorphism recording permutation and orientation is
an isomorphism onto \(\mathfrak S_n\times C_2\) for
\(0\leq n\leq3\), with \(\mathfrak S_0=\mathfrak S_1=1\).
This proves the sphere rows and the assertions about pure groups.

On the torus, action on homology identifies the oriented group
with \(\operatorname{SL}_2(\mathbb Z)\) by the cited
Theorem~2.5; all integral invertible matrices are realized by
linear maps, and orientation is their determinant.
Forgetting one marked point is an isomorphism on mapping-class
groups.  To check injectivity without an erroneous point-pushing
kernel, put that point at \(0\).  If \(f\) fixes \(0\) and
\(F_t\) is an isotopy from \(f\) to the identity on the
unmarked torus, then \(x\mapsto F_t(x)-F_t(0)\) is an
isotopy fixing \(0\).  Surjectivity follows from the linear
representatives, which already fix \(0\).

A reduced-word calculation gives
\(C_{F_2}(a)=\langle a\rangle\) and
\(C_{F_2}(b)=\langle b\rangle\), so \(Z(F_2)=1\).
The other surface-group centres follow from commutativity.
The centre of \(\mathfrak S_3\) is trivial, since commuting
with both \((12)\) and \((23)\) forces a permutation to be
the identity.  A matrix commuting with
\(\left(\begin{smallmatrix}1&1\\0&1\end{smallmatrix}\right)\)
and \(\left(\begin{smallmatrix}1&0\\1&1\end{smallmatrix}\right)\)
is scalar; an invertible scalar integral matrix is \(\pm I\).
This computes the centres of both torus groups and proves the table.
\end{proof}

\noindent\textbf{Scope of the comparisons.}
The first four types have \(2g-2+n\leq0\).
Their geometric groups are trivial or abelian, and their
profinite completions are respectively
\(1,1,\widehat{\mathbb Z},\widehat{\mathbb Z}^{2}\),
with those same groups as centres.  Hence arguments requiring
a centre-free hyperbolic geometric kernel do not apply; no
hyperbolic anabelian reconstruction theorem is being asserted
for these rows.  In particular the forgetful torus map above
has no injected \(\mathbb Z^2\) point-pushing kernel.
The types \((0,3)\) and \((1,1)\) \emph{are} hyperbolic.
Their abstract groups are both \(F_2=\langle a,b\rangle\),
but their peripheral classes are respectively
\(a,b,(ab)^{-1}\) and \([a,b]^{-1}\).
Their distinct peripheral data and mapping-class centres must
therefore be retained in a comparison.

Even for the cylinder, orientation and loop action differ:
if \(\epsilon\) is orientation sign and \(s\) permutation
sign, the induced sign on \(\Pi=\mathbb Z\) is \(s\epsilon\).
Indeed \(z\mapsto1/z\) exchanges the two punctures and
preserves orientation, whereas \(z\mapsto\overline z\)
fixes them and reverses orientation; both act by \(-1\)
on \(\Pi\).  This also explains why isotopy cannot be
replaced here by arbitrary homotopy.

These are discrete topological groups, not the affine
group-scheme automorphism functors in \S\S48--50.
For example the oriented mapping-class group of the sphere
is trivial while its complex algebraic automorphism group
is \(\operatorname{PGL}_2(\mathbb C)\).
Orientation-reversing mapping classes still belong to an
orientable surface; they do not supply a non-orientable version
of the complex-curve comparison.  A real structure adds a
specified antiholomorphic involution and its outer action,
not just a choice of the abstract extended group.
The supplied LM3 A.11 was used as a lead; the calculation and
the cited primary exposition were checked independently here.
